% Source slide 51
\section[4.4 Coordinate Systems]{4.4 Coordinate Systems}
\chaptercard{4}{Vector Spaces}{4.4}{COORDINATE SYSTEMS}

% Source slide 52
\begin{frame}[shrink=7]{THE UNIQUE REPRESENTATION THEOREM}
\begin{itemize}
\item \sourcehead{Theorem 7:} Let $B=\{\vect b_1,\ldots,\vect b_n\}$ be a basis for vector space $V$. Then for each $\vect x$ in $V$, there exists a unique set of scalars $c_1,\ldots,c_n$ such that
\begin{equation}
\vect x=c_1\vect b_1+\cdots+c_n\vect b_n. \tag{1}
\end{equation}
\item \textbf{Proof:} Since $B$ spans $V$, there exist scalars such that (1) holds.
\item Suppose $\vect x$ also has the representation
\[
\vect x=d_1\vect b_1+\cdots+d_n\vect b_n
\]
for scalars $d_1,\ldots,d_n$.
\end{itemize}
\end{frame}

% Source slide 53
\begin{frame}[shrink=9]{THE UNIQUE REPRESENTATION THEOREM}
\begin{itemize}
\item Then, subtracting, we have
\begin{equation}
0=\vect x-\vect x=(c_1-d_1)\vect b_1+\cdots+(c_n-d_n)\vect b_n. \tag{2}
\end{equation}
\item Since $B$ is linearly independent, the weights in (2) must all be zero. That is, $c_j=d_j$ for $1\leq j\leq n$.
\item \textbf{Definition:} Suppose $B=\{\vect b_1,\ldots,\vect b_n\}$ is a basis for $V$ and $\vect x$ is in $V$. The \textbf{coordinates of $\vect x$ relative to the basis $B$} (or the \textbf{$B$-coordinate of $\vect x$}) are the weights $c_1,\ldots,c_n$ such that
\[
\vect x=c_1\vect b_1+\cdots+c_n\vect b_n.
\]
\end{itemize}
\end{frame}

% Source slide 54
\begin{frame}[shrink=7]{THE UNIQUE REPRESENTATION THEOREM}
\begin{itemize}
\item If $c_1,\ldots,c_n$ are the $B$-coordinates of $\vect x$, then the vector in $\R^n$
\[
\coord{\vect x}{B}=\begin{bmatrix}c_1\\ \vdots\\ c_n\end{bmatrix}
\]
is the \textbf{coordinate vector of $\vect x$ (relative to $B$)}, or the \textbf{$B$-coordinatevector of$\vect x$}.
\item The mapping $\vect x\mapsto\coord{\vect x}{B}$ is the \textbf{coordinate mapping (determined by $B$)}.
\end{itemize}
\end{frame}

% Source slide 55
\begin{frame}[shrink=12]{COORDINATES IN $\R^n$}
\begin{itemize}
\item When a basis $B$ for $\R^n$ is fixed, the $B$-coordinate vector of a specified $\vect x$ is easily found, as in the example below.
\item \textbf{Example 1:} Let
\[
\vect b_1=\begin{bmatrix}2\\1\end{bmatrix},\qquad
\vect b_2=\begin{bmatrix}-1\\1\end{bmatrix},\qquad
\vect x=\begin{bmatrix}4\\5\end{bmatrix},
\]
and $B=\{\vect b_1,\vect b_2\}$. Find the coordinate vector $\coord{\vect x}{B}$ of $\vect x$ relative to $B$.
\item \textbf{Solution:} The $B$-coordinate $c_1,c_2$ of $\vect x$ satisfy
\[
c_1\underbrace{\begin{bmatrix}2\\1\end{bmatrix}}_{\vect b_1}
+c_2\underbrace{\begin{bmatrix}-1\\1\end{bmatrix}}_{\vect b_2}
=\underbrace{\begin{bmatrix}4\\5\end{bmatrix}}_{\vect x}.
\]
\end{itemize}
\end{frame}

% Source slide 56
\begin{frame}[shrink=8]{COORDINATES IN $\R^n$}
or
\begin{equation}
\underbrace{\begin{bmatrix}2&-1\\1&1\end{bmatrix}}_{\vect b_1\quad\vect b_2}
\begin{bmatrix}c_1\\c_2\end{bmatrix}
=\underbrace{\begin{bmatrix}4\\5\end{bmatrix}}_{\vect x}. \tag{3}
\end{equation}
\begin{itemize}
\item This equation can be solved by row operations on an augmented matrix or by using the inverse of the matrix on the left.
\item In any case, the solution is $c_1=3$, $c_2=2$.
\item Thus $\vect x=3\vect b_1+2\vect b_2$ and
\[
\coord{\vect x}{B}=\begin{bmatrix}c_1\\c_2\end{bmatrix}
=\begin{bmatrix}3\\2\end{bmatrix}.
\]
\end{itemize}
\end{frame}

% Source slide 57
\begin{frame}[shrink=8]{COORDINATES IN $\R^n$}
\begin{itemize}
\item See the following figure.
\end{itemize}
\begin{columns}[T,onlytextwidth]
\begin{column}{0.48\textwidth}
\centering
\begin{tikzpicture}[scale=.72,>=Stealth]
  \coordinate (O) at (0,0);
  \coordinate (b1) at (1.65,.48);
  \coordinate (b2) at (-.95,.70);
  \foreach \i in {-1,...,4}{
    \draw[gray!65] ($(O)+\i*(b1)-1.2*(b2)$) -- ($(O)+\i*(b1)+3.0*(b2)$);
  }
  \foreach \j in {-1,...,3}{
    \draw[gray!65] ($(O)-1.2*(b1)+\j*(b2)$) -- ($(O)+4.0*(b1)+\j*(b2)$);
  }
  \draw[->] (-3.2,0)--(5.4,0);
  \draw[->] (0,-1.0)--(0,4.0);
  \draw[diagramblue,very thick,->] (O)--(b1) node[below right,black]{$\vect b_1$};
  \draw[diagramblue,very thick,->] (O)--(b2) node[below left,black]{$\vect b_2$};
  \coordinate (x) at ($(O)+3*(b1)+2*(b2)$);
  \fill[diagramblue] (x) circle (2.5pt) node[above right,black]{$\vect x$};
  \fill[diagramblue] (O) circle (2pt);
\end{tikzpicture}

\small The $B$-coordinate vector of $\vect x$\\is $(3,2)$.
\end{column}
\begin{column}{0.50\textwidth}
\begin{itemize}
\item The matrix in (3) changes the $B$-coordinates of a vector $\vect x$ into the standard coordinates for $\vect x$.
\item An analogous change of coordinates can be carried out in $\R^n$ for a basis $B=\{\vect b_1,\ldots,\vect b_n\}$.
\item Let
\[
P_B=\begin{bmatrix}\vect b_1&\vect b_2&\cdots&\vect b_n\end{bmatrix}.
\]
\end{itemize}
\end{column}
\end{columns}
\end{frame}

% Source slide 58
\begin{frame}[shrink=8]{COORDINATES IN $\R^n$}
\begin{itemize}
\item Then the vector equation
\[
\vect x=c_1\vect b_1+c_2\vect b_2+\cdots+c_n\vect b_n
\]
is equivalent to
\begin{equation}
\vect x=P_B\coord{\vect x}{B}. \tag{4}
\end{equation}
\item $P_B$ is called the \textbf{change-of-coordinates matrix from $B$ to the standard basis in $\R^n$}.
\item Left-multiplication by $P_B$ transforms the coordinate vector $\coord{\vect x}{B}$ into $\vect x$.
\item Since the columns of $P_B$ form a basis for $\R^n$, $P_B$ is invertible (by the Invertible Matrix Theorem).
\end{itemize}
\end{frame}

% Source slide 59
\begin{frame}[shrink=7]{COORDINATES IN $\R^n$}
\begin{itemize}
\item Left-multiplication by $P_B^{-1}$ converts $\vect x$ into its $B$-coordinate vector:
\[
P_B^{-1}\vect x=\coord{\vect x}{B}.
\]
\item The correspondence $\vect x\mapsto\coord{\vect x}{B}$, produced by $P_B^{-1}$, is the coordinate mapping.
\item Since $P_B^{-1}$ is an invertible matrix, the coordinate mapping is a one-to-one linear transformation from $\R^n$ onto $\R^n$, by the Invertible Matrix Theorem.
\end{itemize}
\end{frame}

% Source slide 60
\begin{frame}[shrink=8]{THE COORDINATE MAPPING}
\begin{itemize}
\item \sourcehead{Theorem 8:} Let $B=\{\vect b_1,\ldots,\vect b_n\}$ be a basis for a vector space $V$. Then the coordinate mapping $\vect x\mapsto\coord{\vect x}{B}$ is a one-to-one linear transformation from $V$ onto $\R^n$.
\item \textbf{Proof:} Take two typical vectors in $V$, say,
\[
\vect u=c_1\vect b_1+\cdots+c_n\vect b_n,
\qquad
\vect w=d_1\vect b_1+\cdots+d_n\vect b_n.
\]
\item Then, using vector operations,
\[
\vect u+\vect v=(c_1+d_1)\vect b_1+\cdots+(c_n+d_n)\vect b_n.
\]
\end{itemize}
\end{frame}

% Source slide 61
\begin{frame}[shrink=8]{THE COORDINATE MAPPING}
\begin{itemize}
\item It follows that
\[
\coord{\vect u+\vect w}{B}
=\begin{bmatrix}c_1+d_1\\\vdots\\c_n+d_n\end{bmatrix}
=\begin{bmatrix}c_1\\\vdots\\c_n\end{bmatrix}
+\begin{bmatrix}d_1\\\vdots\\d_n\end{bmatrix}
=\coord{\vect u}{B}+\coord{\vect w}{B}.
\]
\item So the coordinate mapping preserves addition.
\item If $r$ is any scalar, then
\[
r\vect u=r(c_1\vect b_1+\cdots+c_n\vect b_n)
=(rc_1)\vect b_1+\cdots+(rc_n)\vect b_n.
\]
\end{itemize}
\end{frame}

% Source slide 62
\begin{frame}[shrink=10]{THE COORDINATE MAPPING}
\begin{itemize}
\item So
\[
\coord{r\vect u}{B}=\begin{bmatrix}rc_1\\\vdots\\rc_n\end{bmatrix}
=r\begin{bmatrix}c_1\\\vdots\\c_n\end{bmatrix}
=r\coord{\vect u}{B}.
\]
\item Thus the coordinate mapping also preserves scalar multiplication and hence is a linear transformation.
\item The linearity of the coordinate mapping extends to linear combinations.
\item If $\vect u_1,\ldots,\vect u_p$ are in $V$ and if $c_1,\ldots,c_p$ are scalars, then
\begin{equation}
\coord{c_1\vect u_1+\cdots+c_p\vect u_p}{B}
=c_1\coord{\vect u_1}{B}+\cdots+c_p\coord{\vect u_p}{B}. \tag{5}
\end{equation}
\end{itemize}
\end{frame}

% Source slide 63
\begin{frame}[shrink=8]{THE COORDINATE MAPPING}
\begin{itemize}
\item In words, (5) says that the $B$-coordinate vector of a linear combination of $\vect u_1,\ldots,\vect u_p$ is the \emph{same} linear combination of their coordinate vectors.
\item The coordinate mapping in Theorem 8 is an important example of an \emph{isomorphism} from $V$ onto $\R^n$.
\item In general, a one-to-one linear transformation from a vector space $V$ onto a vector space $W$ is called an \textbf{isomorphism} from $V$ onto $W$.
\item The notation and terminology for $V$ and $W$ may differ, but the two spaces are indistinguishable as vector spaces.
\end{itemize}
\end{frame}

% Source slide 64
\begin{frame}[shrink=11]{THE COORDINATE MAPPING}
\begin{itemize}
\item \emph{Every vector space calculation in $V$ is accurately reproduced in $W$, and vice versa.}
\item In particular, any real vector space with a basis of $n$ vectors is indistinguishable from $\R^n$.
\item \textbf{Example 7:} Let
\[
\vect v_1=\begin{bmatrix}3\\6\\2\end{bmatrix},\qquad
\vect v_2=\begin{bmatrix}-1\\0\\1\end{bmatrix},\qquad
\vect x=\begin{bmatrix}3\\12\\7\end{bmatrix},
\]
and $B=\{\vect v_1,\vect v_2\}$. Then $B$ is a basis for $H=\Span\{\vect v_1,\vect v_2\}$. Determine if $\vect x$ is in $H$, and if it is, find the coordinate vector of $\vect x$ relative to $B$.
\end{itemize}
\end{frame}

% Source slide 65
\begin{frame}{THE COORDINATE MAPPING}
\begin{itemize}
\item \textbf{Solution:} If $\vect x$ is in $H$, then the following vector equation is consistent:
\[
c_1\begin{bmatrix}3\\6\\2\end{bmatrix}
+c_2\begin{bmatrix}-1\\0\\1\end{bmatrix}
=\begin{bmatrix}3\\12\\7\end{bmatrix}.
\]
\item The scalars $c_1$ and $c_2$, if they exist, are the $B$-coordinates of $\vect x$.
\end{itemize}
\end{frame}

% Source slide 66
\begin{frame}{THE COORDINATE MAPPING}
\begin{itemize}
\item Using row operations, we obtain
\[
\begin{bmatrix}3&-1&3\\6&0&12\\2&1&7\end{bmatrix}
\sim
\begin{bmatrix}1&0&2\\0&1&3\\0&0&0\end{bmatrix}.
\]
\item Thus $c_1=2$, $c_2=3$ and
\[
\coord{\vect x}{B}=\begin{bmatrix}2\\3\end{bmatrix}.
\]
\end{itemize}
\end{frame}

% Source slide 67
\begin{frame}{THE COORDINATE MAPPING}
\begin{itemize}
\item The coordinate system on $H$ determined by $B$ is shown in the following figure.
\end{itemize}
\centering
\begin{tikzpicture}[scale=.72,>=Stealth]
  \coordinate (O) at (0,0);
  \coordinate (v1) at (2.2,-.25);
  \coordinate (v2) at (.72,1.42);
  \fill[blue!18] ($(O)-.25*(v1)-.25*(v2)$) -- ($(O)+2.75*(v1)-.25*(v2)$) -- ($(O)+2.75*(v1)+3.1*(v2)$) -- ($(O)-.25*(v1)+3.1*(v2)$) -- cycle;
  \foreach \i in {0,...,3}{\draw[black!75] ($(O)+\i*(v1)$)--($(O)+\i*(v1)+3*(v2)$);}
  \foreach \j in {0,...,3}{\draw[black!75] ($(O)+\j*(v2)$)--($(O)+2.6*(v1)+\j*(v2)$);}
  \draw[diagramblue,very thick,->] (O)--($(O)+2.75*(v1)$);
  \draw[diagramblue,very thick,->] (O)--($(O)+3.15*(v2)$);
  \draw[->] (-2.2,-1.9)--(6.8,2.9);
  \draw[->] (0,-.8)--(0,4.9);
  \fill[diagramblue] (v1) circle (2.5pt) node[below]{$\vect v_1$};
  \fill[diagramblue] ($(O)+2*(v1)$) circle (2.5pt) node[below]{$2\vect v_1$};
  \fill[diagramblue] (v2) circle (2.5pt) node[left]{$\vect v_2$};
  \fill[diagramblue] ($(O)+2*(v2)$) circle (2.5pt) node[left]{$2\vect v_2$};
  \fill[diagramblue] ($(O)+3*(v2)$) circle (2.5pt) node[left]{$3\vect v_2$};
  \coordinate (x) at ($(O)+2*(v1)+3*(v2)$);
  \fill[diagramblue] (x) circle (2.8pt) node[right,black]{$\vect x=2\vect v_1+3\vect v_2$};
  \draw[dashed] (x)--($(O)+2*(v1)$);
  \draw[dashed] (x)--($(O)+3*(v2)$);
  \node[below left] at (O) {$0$};
\end{tikzpicture}

\small A coordinate system on a plane $H$ in $\R^3$.
\end{frame}
